% !TEX TS-program = pdflatex
% !TeX program = pdflatex
% !TEX encoding = UTF-8
% !TEX spellcheck = en_US

\documentclass[xcolor=table]{beamer}

\usepackage{../extra/beamer/karimnlp}

\input{options}

\subtitle[01- Introduction]{Chapter 01\\Introduction to NLP} 

\changegraphpath{../img/intro/}

\begin{document}
	
\begin{frame}
	\frametitle{\inserttitle}
	\framesubtitle{\insertshortsubtitle: Definitions}
	
	\begin{itemize}
		\item \textbf{NLP (Natural Language Processing)}: the study and application of computational techniques to process, understand, and generate human language.
	\end{itemize}
	\begin{minipage}{0.78\textwidth}
		\begin{itemize}
			\item A subfield of:
			\begin{itemize}
				\item \optword{Linguistics}: the scientific study of human language.
				\item \optword{Computer Science}: the study of computation, automated systems, and information processing.
				\item \optword{Artificial Intelligence}: the science and engineering of designing intelligent systems \cite{1955-mccarthy}.
			\end{itemize}
		\end{itemize}
	\end{minipage}
	\begin{minipage}{0.20\textwidth}
		\hgraphpage[\textwidth]{NLP.pdf}
	\end{minipage}

\end{frame}

\begin{frame}
	\frametitle{\inserttitle}
	\framesubtitle{\insertshortsubtitle: Motivation}
	
	\begin{itemize}
		\item The sheer volume of textual data generated daily.
		\item Most textual data is unstructured, making analysis challenging.
		\item Textual communication spans multiple domains (social media, email, forums, etc.).
		\item Growing adoption of NLP techniques in industry and research.
		\item Increasing interest in applications such as chatbots, sentiment analysis, and information retrieval.
	\end{itemize}
	
	\begin{center}
		{\Large NLP has become a central focus in research and industry.}
	\end{center}

\end{frame}

\begin{frame}
	\frametitle{\inserttitle}
	\framesubtitle{\insertshortsubtitle: Some humor}
	
	\begin{center}
		\vgraphpage{humor/humor.jpg}
	\end{center}

\end{frame}

\begin{frame}
	\frametitle{\inserttitle}
	\framesubtitle{\insertshortsubtitle: Plan}
	
	\begin{multicols}{2}
		%	\small
		\tableofcontents
	\end{multicols}

\end{frame}

%===================================================================================
\section{History}
%===================================================================================

\begin{frame}
	\frametitle{\insertshortsubtitle}
	\framesubtitle{\insertsection}

	\hgraphpage{history.pdf}

\end{frame}

\subsection{AI birth and golden era}

\begin{frame}
	\frametitle{\insertshortsubtitle: \insertsection}
	\framesubtitle{\insertsubsection: The 1950s}

	\begin{itemize}
		\item \optword{1951}: Shannon analyzed probabilistic models of language, laying foundations for information-theoretic approaches \cite{1951-shannon}.
		\item \optword{1954}: Georgetown-IBM experiment automatically translated 60 Russian sentences into English.
		\item \optword{1956}: Chomsky proposed a hierarchy of formal grammars, known as the \keyword{Chomsky hierarchy}.
		\item \optword{1958}: Luhn experimented with automatic text summarization via extraction techniques \cite{1958-luhn}.
	\end{itemize}

\end{frame}

\begin{frame}
	\frametitle{\insertshortsubtitle: \insertsection}
	\framesubtitle{\insertsubsection: The 1960s}

	\begin{itemize}
		\item \optword{1961}: Development of the first automatic English parser at UPenn \cite{1961-joshi,1962-harris}.
		\item \optword{1964}: Weizenbaum developed \keyword{ELIZA}, a simulation of a psychotherapist at MIT AI Lab.
		\item \optword{1964}: Bobrow developed \keyword{STUDENT}, designed to read and solve algebra word problems \cite{1964-bobrow}.
		\item \optword{1967}: Creation of the \keyword{Brown Corpus}, the first electronic corpus of English text.
	\end{itemize}

\end{frame}

\subsection{AI winter}

\begin{frame}
	\frametitle{\insertshortsubtitle: \insertsection}
	\framesubtitle{\insertsubsection: The 1970s}

	\begin{itemize}
		\item \optword{1971}: Winograd developed \keyword{SHRDLU}, a natural language understanding program \cite{1971-winograd}.
		\item \optword{1972}: Colby created \keyword{PARRY}, simulating a person with paranoid schizophrenia.
		\item \optword{1975}: \keyword{MARGIE} performed inference and paraphrasing using conceptual representations of language.
		\item \optword{1975}: \keyword{DRAGON}, an automatic speech recognition system using hidden Markov models \cite{1975-baker}.
	\end{itemize}

\end{frame}

\begin{frame}
	\frametitle{\insertshortsubtitle: \insertsection}
	\framesubtitle{\insertsubsection: The 1980s}

	\begin{itemize}
		\item \optword{1980}: \keyword{KL-One}, knowledge representation system for syntax and semantics \cite{1980-bobrow}.
		\item \optword{1986}: \keyword{TRUMP}, language analyzer using a lexical database \cite{1986-jacobs}.
		\item \optword{1987}: \keyword{HPSG} (Head-driven Phrase Structure Grammar) \cite{1987-sag-pollard}.
		\item \optword{1987}: \keyword{MUC} (Message Understanding Conference), sponsored by \keyword{DARPA}.
		\item \optword{1988}: Application of hidden Markov models for part-of-speech tagging \cite{1988-church}.
		\item Symbolic methods for speech processing and natural language generation.
	\end{itemize}

\end{frame}

\subsection{AI spring}

\begin{frame}
	\frametitle{\insertshortsubtitle: \insertsection}
	\framesubtitle{\insertsubsection: The 1990s}
	
	\begin{itemize}
		\item \optword{1990}: Statistical approaches to machine translation \cite{1990-brown-al}.
		\item \optword{1993}: \keyword{Penn Treebank}, an annotated English corpus \cite{1993-marcus-al}.
		\item \optword{1995}: \keyword{Wordnet}, a lexical database of English \cite{1995-miller}.
		\item \optword{1996}: \keyword{SPATTER}, a statistical lexical analyzer using decision trees \cite{1996-magerman}.
		\item Trend: Widespread adoption of statistical methods and empirical evaluation.
	\end{itemize}

\end{frame}

\begin{frame}
	\frametitle{\insertshortsubtitle: \insertsection}
	\framesubtitle{\insertsubsection: The 2000s}
	
	\begin{itemize}
		\item \optword{2003}: Introduction of neural language models \cite{2003-bengio-al}.
		\item \optword{2006}: IBM developed \keyword{Watson}, a question-answering system.
		\item Trend: Increased use of unsupervised and semi-supervised learning.
		\item Trend: Shift toward semantic-level NLP tasks.
	\end{itemize}

\end{frame}

\begin{frame}
	\frametitle{\insertshortsubtitle: \insertsection}
	\framesubtitle{\insertsubsection: The 2010s}
	
	\begin{itemize}
		\item \optword{2011}: \keyword{Siri} (Apple), a personal digital assistant. Followed by \keyword{Alexa} (Amazon, 2014) and \keyword{Google Assistant} (2016).
		\item \optword{2014}: Introduction of word embeddings \cite{2014-lebret-collobert}.
		\item \optword{2018}: Emergence of contextualized pre-trained language models:
		\begin{itemize}
			\item \keyword{ULMfit} (fast.ai) \cite{2018-howard-ruder}
			\item \keyword{ELMO} (AllenNLP) \cite{2018-peters-al}
			\item \keyword{GPT} (OpenAI) \cite{2018-radford-al}
			\item \keyword{BERT} (Google) \cite{2018-devlin-al}
			\item \keyword{XLM} (Facebook, 2019) \cite{2019-lample-conneau}
		\end{itemize}
	\end{itemize}

\end{frame}

\begin{frame}
	\frametitle{\insertshortsubtitle: \insertsection}
	\framesubtitle{\insertsubsection: The 2020s}
	
	\begin{itemize}
		\item \optword{2020}: Focus on interpretability and explainability of neural NLP models \cite{2020-belinkov-al,2020-kim-al,2020-jacovi-goldberg}.
		\item \optword{2021}: Research on ethics and fairness in NLP \cite{2021-garridoMunoz-al,2021-field-al}.
		\item \optword{2022}: Development of polyglot language models such as \keyword{Polyglot Prompt} \cite{2022-fu-al}.
		\item \optword{2022}: \keyword{Edge NLP}: deploying efficient NLP models on resource-constrained devices \cite{2022-guo-al}.
	\end{itemize}
	
\end{frame}

%===================================================================================
\section{Language processing levels}
%===================================================================================

\begin{frame}
	\frametitle{\insertshortsubtitle}
	\framesubtitle{\insertsection}

	\begin{center}
		\vgraphpage{levels.pdf}
	\end{center}

\end{frame}

\subsection{Phonetics, phonology \& orthography}

\begin{frame}
	\frametitle{\insertshortsubtitle: \insertsection}
	\framesubtitle{\insertsubsection: Phonetics}

	\begin{itemize}
		\item Study of the sounds (phones) produced by the human vocal apparatus.
		\item Branches:
		\begin{itemize}
			\item \optword{Articulatory phonetics}: anatomical and physiological branch; describes how vowels and consonants are produced in various parts of the mouth and throat.
			\item \optword{Acoustic phonetics}: branch closely associated with physics; studies the sound waves that transmit vowels and consonants from speaker to listener.
			\item \optword{Auditory phonetics}: branch concerned with perception; examines how listeners interpret sound waves to perceive speech sounds.
		\end{itemize}
	\end{itemize}

\end{frame}

\begin{frame}
	\frametitle{\insertshortsubtitle: \insertsection}
	\framesubtitle{\insertsubsection: Phonetics (Articulation points of consonants)}

	\begin{minipage}{0.55\textwidth}
	\begin{itemize}
		\item \optword{labial}: Lips. E.g., \expword{\textipa{[b], [p], [m], [f], [v]}}.
		\item \optword{apical}: Tip or front of the tongue. 
		E.g., \expword{%
			\textipa{[t], [d], [n], [r], }
			\RL{^s} \textipa{[S],} 
			\RL{_t} \textipa{[T],} 
			\RL{_d} \textipa{[D]}
		}
		\item \optword{dorsal}: Back of the tongue. E.g., \expword{\textipa{[c], [k], [g], [q],} \RL{.g} \textipa{[G]}}.
		\item \optword{pharyngeal}: Pharynx. 
		E.g., \expword{\RL{.h} \textipa{[\*h],} \RL{`} \textipa{[Q]}}.
		\item \optword{glottal}: Glottis. 
		E.g., \expword{\textipa{[h],} \RL{'} \textipa{[P]}}.
	\end{itemize}
	\end{minipage}
	\begin{minipage}{0.43\textwidth}
		\begin{figure}
			\hgraphpage{vocaltract.pdf}\caption{Vocal tract \cite{2009-ball}}
		\end{figure}
	\end{minipage}

\end{frame}

\begin{frame}
	\frametitle{\insertshortsubtitle: \insertsection}
	\framesubtitle{\insertsubsection: Phonetics (Articulation modes of consonants)}

	\begin{itemize}
		\item \optword{occlusive}: complete blockage of airflow.
		E.g., \expword{\textipa{[p], [k], [b], [m], [n]}}.
		
		\item \optword{fricative}: mouth, pharynx or glottis tightening without their complete closure.
		E.g., \expword{\textipa{[f], [v], [s]}}.
		
		\item \optword{Affricate}: airflow is blocked then released.
		E.g., \expword{\textipa{[\t{\textteshlig}]}}.
		
		\item \optword{lateral}: airflow through a lateral channel (sometimes bilateral).
		E.g., \expword{\textipa{[l]}}.
		
		\item \optword{nasal}: lowering the soft palate.
		E.g., \expword{\textipa{[m], [n]}}.
		
		\item \optword{click}: with the tongue or lips without using the lungs.
		E.g., \expword{\textipa{[!]} (tongue clicking)}.
		
	\end{itemize}

\end{frame}

\begin{frame}
	\frametitle{\insertshortsubtitle: \insertsection}
	\framesubtitle{\insertsubsection: Phonetics (IPA)}

	\begin{center}
		\vgraphpage{IPA2020_.pdf}
	\end{center}

\end{frame}

\begin{frame}
	\frametitle{\insertshortsubtitle: \insertsection}
	\framesubtitle{\insertsubsection: Phonology}

	\begin{itemize}
		\item Study of the sounds (phonemes) of a language.
		\item Focuses on how sounds function as elements of a system.
	\end{itemize}
	
	\begin{exampleblock}{Example: Phoneme \textit{/r/}}
		\begin{itemize}
			\item \textbf{French}: /r/ has multiple phonetic realizations:
			\begin{itemize}
				\item Trill (rolled R): \expword{\textipa{[r]}}
				\item Voiced uvular trill: \expword{\textipa{[\;R]}}
				\item Voiced uvular fricative (Parisian): \expword{\textipa{[K]}}
			\end{itemize}
			\item Despite these variations, the phoneme is consistently transcribed as: \expword{/r/}, e.g., \textit{rat /rat/}.
			
			\item In Arabic, consonants \expword{\RL{r} \textipa{[r]}} and \expword{\RL{.g} \textipa{[G][K]}} represent different phonemes: \expword{\textipa{/r/}} and \expword{\textipa{/G/ /K/}} respectively. 
			E.g., \expword{\RL{.gryb} \textipa{/G\ae ri:b/} (stranger)}.
		\end{itemize}
	\end{exampleblock}

\end{frame}

\begin{frame}
	\frametitle{\insertshortsubtitle: \insertsection}
	\framesubtitle{\insertsubsection: Orthography}

	\begin{itemize}
		\item Study of writing systems, graphemes, and spelling conventions.
		\item Writing systems (based on \keyword{graphemes}):
		\begin{itemize}
			\item \optword{Logographic}: Each grapheme represents a word or morpheme. 
			E.g., \expword{Kanji (Japanese): \begin{CJK}{UTF8}{min}日, 本, 語\end{CJK}}.
			
			\item \optword{Syllabic}: Each symbol represents a syllable. 
			E.g., \expword{Hiragana (Japanese): \begin{CJK}{UTF8}{min}る, た, め\end{CJK}; Katakana: \begin{CJK}{UTF8}{min}セ, ク\end{CJK}}.
			
			\item \optword{Alphabetic}: Each letter represents a phoneme. 
			E.g., \expword{Latin: A, B, C, ...}.
			
			\item \optword{Abjad}: Similar to alphabetic systems, but primarily consonants are represented. 
			E.g., \expword{Arabic: \RL{b}, \RL{t}, \RL{h}}.
		\end{itemize}
		\item Punctuation symbols.
		\item Writing rules and orthographic conventions.
	\end{itemize}

\end{frame}

\begin{frame}
	\frametitle{\insertshortsubtitle: \insertsection}
	\framesubtitle{\insertsubsection: Some humor}

	\begin{center}
		\vgraphpage{humor/humor-phonetics.jpg}
	\end{center}

\end{frame}

\subsection{Morphology and syntax}

\begin{frame}
	\frametitle{\insertshortsubtitle: \insertsection}
	\framesubtitle{\insertsubsection: Morphology}
	
	\begin{itemize}
		\item Study of word formation, including word creation and derivation.
		\item Study of how words change form according to grammatical context.
		\item \keyword{Morpheme}: the smallest meaningful unit of language. 
		E.g., \expword{proper nouns, suffixes, prefixes, etc.}
		\item \keyword{Lexeme}: an abstract unit representing all inflected forms of a word.
		\item \keyword{Lemma}: a canonical form chosen to represent a lexeme.
		\item Parts of speech:
		\begin{itemize}
			\item Open class: \optword{adjective}, \optword{noun}, \optword{verb}.
			\item Closed class: \optword{adverb}, \optword{article}, \optword{conjunction}, \optword{interjection}, \optword{preposition}, \optword{pronoun}.
		\end{itemize}
	\end{itemize}

\end{frame}


\begin{frame}
	\frametitle{\insertshortsubtitle: \insertsection}
	\framesubtitle{\insertsubsection: Morphology}
%
%	\begin{itemize}
%		\item \optword{Isolating/analytical languages}: each word consists of one and only one morpheme. 
%		Morphological changes are few or even absent.
%		Among these languages: \expword{Mandarin, Vietnamese, Tai, Khmer, etc.}. 
%		E.g., \expword{\begin{CJK}{UTF8}{min}四个男孩\end{CJK} /sì ge nánhái/ ``four boys" (lit. ``four [entity of] masculine child")}.
%		\item \optword{Inflectional/synthetic languages}: words are formed from a \keyword{root} plus additional morphemes.
%		\begin{itemize}
%			\item \optword{Agglutinating languages}: morphemes are clearly phonetically distinguishable from each other.
%			Among these languages: \expword{Finnish, Turkish, Japanese, etc.}. 
%			E.g., \expword{\begin{CJK}{UTF8}{min}行く, 行きます\end{CJK}}.
%			\item \optword{Fusional languages}: it is not always easy to distinguish morphemes from root, or from each other.
%			Among these languages: \expword{Arabic, English, French, Berber languages, etc.}.
%			E.g., \expword{\RL{kitAb, kutub, 'wa'`.taynAkumuwh}, foot, feet}
%		\end{itemize}
%	\end{itemize}
	
	\begin{itemize}
		\item \optword{Isolating/Analytical languages}: each word consists of a single morpheme; minimal inflection.
		\begin{itemize}
			\item Examples: \expword{Mandarin, Vietnamese, Tai, Khmer}.
			\item E.g., \expword{\begin{CJK}{UTF8}{min}四个男孩\end{CJK} /sì ge nánhái/ ``four boys'' (lit. ``four [entity of] masculine child'')}.
		\end{itemize}
		
		\item \optword{Inflectional/Synthetic languages}: words are formed from a root plus additional morphemes.
		\begin{itemize}
			\item \optword{Agglutinating languages}: morphemes are clearly separable.
			\begin{itemize}
				\item Examples: \expword{Finnish, Turkish, Japanese}.
				\item E.g., \expword{\begin{CJK}{UTF8}{min}行く, 行きます\end{CJK}}.
			\end{itemize}
			
			\item \optword{Fusional languages}: morphemes are not always easily separable from the root.
			\begin{itemize}
				\item Examples: \expword{Arabic, English, French, Berber}.
				\item E.g., English: \expword{foot → feet}; Arabic: \expword{\RL{kitAb}  →  \RL{kutub}}.
			\end{itemize}
		\end{itemize}
	\end{itemize}

\end{frame}

\begin{frame}
	\frametitle{\insertshortsubtitle: \insertsection}
	\framesubtitle{\insertsubsection: Inflectional morphology}
	
	\begin{itemize}
		\item Formation of word variants without changing lexical category or creating new lexemes.
		\item Modifying the form of lexemes to adapt to different grammatical contexts.
		\item \optword{Flexion}: \optword{Declension} (nouns, adjectives) or \optword{Conjugation} (verbs).
		\item \optword{Affixation}: adding morphemes to a root to indicate grammatical features.
		\begin{itemize}
			\item \optword{Prefix}: \expword{\RL{_dhb}} (past), \expword{\RL{y_dhb}} (present).
			\item \optword{Infix}: e.g., Tagalog \expword{sulat → sumulat}.
			\item \optword{Suffix}: \expword{étudiant (singular-masculine)}, \expword{étudiantes (plural-feminine)}.
		\end{itemize}
		\item \optword{Reduplication}: repetition of full or part of a word to express grammatical or semantic meaning. 
		E.g., \expword{super-duper, bye-bye, \begin{CJK}{UTF8}{min}きらきら\end{CJK}}.
	\end{itemize}

\end{frame}

\begin{frame}
	\frametitle{\insertshortsubtitle: \insertsection}
	\framesubtitle{\insertsubsection: Inflectional morphology (Grammatical features)}

	\begin{itemize}
		\item \optword{Number}: singular, dual, trial, paucal, plural, etc. 
		\item \optword{Person}: first, second, third, etc.
		\item \optword{Gender}: masculine, feminine, neuter, common.
		\item \optword{Case}: nominative, absolutive, accusative, ergative, etc.
		\item \optword{Tense}: past, present, future.
		\item \optword{Aspect}: accomplished/unaccomplished, progressive, perfective/imperfective, iterative, prospective.
		\item \optword{Voice}: active, medium, passive, etc.
		\item \optword{Polarity}: affirmative, negative.
		\item \optword{Politeness}: informal, formal, etc.
	\end{itemize}

\end{frame}

\begin{frame}
	\frametitle{\insertshortsubtitle: \insertsection}
	\framesubtitle{\insertsubsection: Derivational morphology}
	\small\vspace{-3pt}
	\begin{itemize}
		\item Formation of new words by changing part of speech (\expword{play, player}) or creating new lexemes (\expword{do, undo}). 
		\item \optword{Affixation}: using prefixes, infixes and/or suffixes. 
		E.g., \expword{happy (ADJ), unhappy (ADJ), unhappiness (N)}; \expword{\RL{jhd} (V), \RL{ijthd} (V)}.
		
		\item \optword{Compounding}: combining two words into a single lexical item.
		E.g., \expword{wind (N) + mill (N) = windmill (N)}.
		
		\item \optword{Blending}: combining fragments of words into a new word.
		E.g., \expword{porter (V) + manteau (N) = portemanteau (N)}; \expword{breakfast (N) + lunch (N) = brunch (N)}
		
		\item \optword{Conversion}: changing a word's part of speech without morphological modification.
		E.g., \expword{orange (fruit, N), orange (color, ADJ)}; \expword{visit-er (V), visite (N)}; \expword{fish (N), to fish (V)}.
		
		\item \optword{Truncation}: shortening a word.
		E.g., \expword{bibliographie, biblio}, \expword{information, info}.
		
		\item \optword{Reduplication}: repetition of a word or part of it.
		E.g., \expword{\RL{kr} (V), \RL{krkr} (V)}.
		
	\end{itemize}

\end{frame}

\begin{frame}
	\frametitle{\insertshortsubtitle: \insertsection}
	\framesubtitle{\insertsubsection: Syntax}

	\begin{itemize}
		\item Sentence structure: how words are combined to form sentences.
		\item Parts of speech: the grammatical category of each word in a sentence (verb, noun, adjective, etc.), relevant for POS tagging.
		\item Word order: typical sequence of constituents, e.g., \keyword{Subject (S)}, \keyword{Verb (V)} and \keyword{Object (O)}. 
		\item Grammatical functions: syntactic roles such as  \expword{Subject, Direct object, Indirect object, etc.}
	\end{itemize}

\end{frame}

\begin{frame}
	\frametitle{\insertshortsubtitle: \insertsection}
	\framesubtitle{\insertsubsection: Syntax (Word order: Example)}
	
	\begin{itemize}
		\item \textcolor{blue}{Subject}, \textcolor{red}{Verb}, \textcolor{green}{Object}
	\end{itemize}
	
	\begin{columns}
		\begin{column}{0.52\textwidth}
			\begin{exampleblock}{Example SOV [Japanese]}
				\begin{CJK}{UTF8}{min}
					\textcolor{blue}{カリムさん}は\textcolor{green}{日本語}を\textcolor{red}{勉強します}。
				\end{CJK}
			\end{exampleblock}
		\end{column}
		%
		\begin{column}{0.42\textwidth}
			\begin{exampleblock}{Example SVO [English]}
				\textcolor{blue}{Karim} \textcolor{red}{learns} \textcolor{green}{English}.
			\end{exampleblock}
		\end{column}%
	\end{columns}
	\begin{columns}
		\begin{column}{0.49\textwidth}
			\begin{exampleblock}{Example VSO [Arabic]}
				\begin{flushright}
					\color{black}.
					\color{green} \RL{al-`rbyT}%
					\color{blue} \RL{krym} 
					\color{red} \RL{yt`llm} 
				\end{flushright}
			\end{exampleblock}
		\end{column}
	\end{columns}
	
\end{frame}

\begin{frame}
	\frametitle{\insertshortsubtitle: \insertsection}
	\framesubtitle{\insertsubsection: Syntax (Word order: Statistics)}

	\begin{table}
		\begin{tblr}{
				colspec = {p{.1\textwidth}p{.15\textwidth}p{.65\textwidth}},
				row{odd} = {lightblue},
				row{even} = {lightyellow},
				row{1} = {darkblue},
			} 
			\textcolor{white}{Order} & \textcolor{white}{Proportion} & \textcolor{white}{Examples} \\
			SOV & 44.78\% & Japanese, Latin, Tamil, Basque, Urdu, Ancient Greek, Bengali, Hindi, Sanskrit, Persian, Korean \\
			SVO & 41.79\% & French, Mandarin, Russian, English, Hausa, Italian, Malay (language), Spanish, Thai\\
			VSO & 9.20\% & Irish, Arabic, Biblical Hebrew, Filipino, Tuareg languages, Welsh \\
			VOS & 2.99\% & Malagasy, Baure, Car (language) \\
			OVS & 1.24 \% & Apalai, Hixkaryana, Klingon (language) \\
		\end{tblr}

		\caption{Words Order proportions based on 402 languages \cite{1988-blake}}
	\end{table}

\end{frame}

\begin{frame}
	\frametitle{\insertshortsubtitle: \insertsection}
	\framesubtitle{\insertsubsection: Syntax (Constituency grammar)}

	\begin{itemize}
		\item A \keyword{sentence} is composed of several \keyword{phrases}, which include words and/or other phrases.
		\item A phrase contains a \keyword{head} (\keyword{nucleus}), the central element determining its syntactic type.
		\item A phrase is classified according to its head: nominal (\keyword{NP}), adjectival (\keyword{AP}), verbal (\keyword{VP}), or prepositional (\keyword{PP}). 
		\item The most widely used formal system to model constituency structure is the \keyword{context-free grammar} (Level 2 in \keyword{Chomsky}'s hierarchy), capable of representing hierarchical sentence structures.
	\end{itemize}

\end{frame}


\begin{frame}
	\frametitle{\insertshortsubtitle: \insertsection}
	\framesubtitle{\insertsubsection: Syntax (Constituency grammar: Example)}

	\begin{columns}
		\begin{column}{0.52\textwidth}
		The grammar that generated this syntactic tree can be written as: 
		
		\begin{enumerate}
			\item S $ \rightarrow $ NP VP
			\item NP $ \rightarrow $ DET NP'
			\item NP $ \rightarrow $ DET N
			\item NP' $ \rightarrow $ ADJ N
			\item VP $ \rightarrow $ V NP
		\end{enumerate}
	
		The second rule is not written as \\(NP $ \rightarrow $ DET NP)\\ 
		otherwise there can be several determiners to a noun.
		
		\end{column}
		%
		\begin{column}{0.42\textwidth}
			\hgraphpage{const_gram.pdf}
		\end{column}%
	\end{columns}

\end{frame}

\begin{frame}
	\frametitle{\insertshortsubtitle: \insertsection}
	\framesubtitle{\insertsubsection: Syntax (Dependency grammar)}
	
	\begin{itemize}
		\item Dependency grammar represents syntactic structure in terms of words rather than phrases.
		\item \optword{Dependency grammar}: a set of binary directed relations between words, where one word (the head) governs another (the dependent).
		\item Relations include, for example: nominal subject (\optword{nsubj}), object (\optword{obj}), adjective modifier (\optword{amod}), determiner (\optword{det}), etc.
	\end{itemize}
	
	\begin{figure}
		\centering
		\hgraphpage[0.6\textwidth]{dep_gram_.pdf}
		\caption{Dependency tree example, generated using \url{https://corenlp.run/}}
	\end{figure}

\end{frame}

\begin{frame}
	\frametitle{\insertshortsubtitle: \insertsection}
	\framesubtitle{\insertsubsection: Some humor}

	\begin{center}
		\vgraphpage{humor/humor-syntax.jpg}
	\end{center}

\end{frame}

\subsection{Semantics}

\begin{frame}
	\frametitle{\insertshortsubtitle: \insertsection}
	\framesubtitle{\insertsubsection}
	
	\begin{itemize}
		\item Study of meaning.
		\begin{itemize}
			\item \optword{Lexical semantics}: the meaning of words.
			\item \optword{Propositional semantics}: the meaning of sentences.
		\end{itemize}
	\end{itemize}
	
	\begin{exampleblock}{Example of lexical ambiguity}
		The word \expword{chicken} has multiple meanings (from \url{http://wordnetweb.princeton.edu/perl/webwn})
		\begin{itemize}
			\item The cat is chasing some \expword{chickens}. ``\textit{a domestic fowl bred for flesh or eggs.}''
			\item I ate some \expword{chicken}. ``\textit{the flesh of a chicken used for food.}''
			\item Don't be a \expword{chicken}! ``\textit{a person who lacks confidence; irresolute and wishy-washy.}''
		\end{itemize}
	\end{exampleblock}

\end{frame}

\begin{frame}
	\frametitle{\insertshortsubtitle: \insertsection}
	\framesubtitle{\insertsubsection: Word meaning}
	
	\begin{itemize}
		\item In logographic systems, a \keyword{grapheme} can represent one or more meanings, sometimes with multiple pronunciations.
		\begin{itemize}
			\item Japanese Kanji examples:
			\begin{CJK}{UTF8}{min}
				\expword{川} (river), \expword{山} (mountain) \\
				\expword{音} (sound, noise) + \expword{楽} (music, comfort, ease) = \expword{音楽} (music)
			\end{CJK}
		\end{itemize}
		
		\item A word can have several meanings (\keyword{Polysemy}).
		
		\item The meaning of a word is constructed:
		\begin{itemize}
			\item from a set of ``semantic primitives''.
			\item through relations between morphemes, forming a network of meanings.
		\end{itemize}
	\end{itemize}

\end{frame}


\begin{frame}
	\frametitle{\insertshortsubtitle: \insertsection}
	\framesubtitle{\insertsubsection: Word meaning (Semic analysis)}

	\begin{itemize}
		\item Proposed by Eric Buyssens in the 1930s and developed by Bernard Pottier in the 1960s.
		\item \keyword{Seme}: a minimal unit of meaning (minimal semantic feature).
		\item \keyword{Sememe}: a cluster of semes corresponding to a lexical unit.
		\item Words that share one or more positive semes belong to the same semantic field.
	\end{itemize}
	
	\begin{exampleblock}{Example of semic analysis}
		\centering
		\begin{tabular}{|l|l|l|l|}
			\hline
			Word/Seme & animate & domestic & feline \\
			\hline
			Cat & + & + & + \\
			\hline
			Lion & + & - & + \\
			\hline
			Dog & + & + & - \\
			\hline
		\end{tabular}
	\end{exampleblock}

\end{frame}

\begin{frame}
	\frametitle{\insertshortsubtitle: \insertsection}
	\framesubtitle{\insertsubsection: Word meaning (Semantic relations)}
	
	\begin{itemize}
		\item \optword{Synonymy}: two words are synonyms if one can be substituted for the other in a sentence without changing its meaning.
		\item \optword{Antonymy}: words with opposite meanings. They usually express two values of the same property.
		E.g., \expword{big} and \expword{small} express the property \expword{size}.
		
		\item \optword{Hyponymy}: a word with a more specific meaning (a more specific term) than another.  
		E.g., \expword{cat} is a hyponym of \expword{feline}, which is a hyponym of \expword{animal}.
		
		\item \optword{Hypernymy}: a word with a more general meaning (a more general term).  
		E.g., \expword{feline} is a hypernym of \expword{cat}.
		
		\item \optword{Meronymy}: a word that denotes a part of another entity.  
		E.g., \expword{wheel} is a meronym of \expword{car}.
	\end{itemize}

\end{frame}

\begin{frame}
	\frametitle{\insertshortsubtitle: \insertsection}
	\framesubtitle{\insertsubsection: Propositional meaning}

	\begin{minipage}{0.5\textwidth}
		\begin{itemize}
			\item First-order logic.
			\item \keyword{Predicate}: denotes a property or a relation over one or more entities. 
			E.g., \expword{Person(x), Possess(x, y)}.	
			\item \keyword{Constant}: denotes a specific entity in the world.
			E.g., \expword{Karim, Algeria, ESI}.
			\item \keyword{Variable}: denotes an unspecified entity. 
			E.g., \expword{x, y}.
			\item \keyword{Function}: a mapping from entities to another entity.
			E.g., \expword{locationOf(ESI)}.
		\end{itemize}
	\end{minipage}
	\begin{minipage}{0.48\textwidth}
		\begin{figure}
			{\tiny\bfseries
			\begin{tabular}{rcl}
				\hline\hline
				Formula & \textrightarrow & AtomicFormula \\
				& \textbar        & Formula Connective Formula \\
				& \textbar        & Quantifier Variable, ... Formula \\
				& \textbar        & $\textlnot$ Formula \\
				& \textbar        & (Formula) \\
				AtomicFormula & \textrightarrow & Predicate (Term, ...) \\
				Term    & \textrightarrow & Function(Term, ...) \\
				& \textbar        & Constant \\
				& \textbar        & Variable \\
				Connective & \textrightarrow & $\wedge$ \textbar $\vee$ \textbar $\Rightarrow$ \\
				Quantifier & \textrightarrow & $\forall$ \textbar $\exists$ \\
				Constant & \textrightarrow & ESI \textbar Karim \textbar Algérie ...\\
				Variable & \textrightarrow & x \textbar y \textbar ... \\
				Predicate & \textrightarrow & Ecole \textbar EnseignatA \textbar Utiliser \textbar ... \\
				Function & \textrightarrow & EmplacementDe \textbar ... \\
				\hline\hline
			\end{tabular}
		}
			\caption{Syntax adapted from \cite{2002-russell-norvig} describing the representation of first-order logic \cite{2019-jurafsky-martin}}
		\end{figure}
	\end{minipage}

\end{frame}

\begin{frame}
	\frametitle{\insertshortsubtitle: \insertsection}
	\framesubtitle{\insertsubsection: Propositional meaning (Example)}

	\begin{exampleblock}{Example of a semantic representation}
		\begin{itemize}
			%	\item I have a car \cite{2019-jurafsky-martin}
			%	\[%
			%	\exists\, e, y\, Having(e) \wedge Haver(e, Speaker) \wedge HadThing(e, y) \wedge Car(y)
			%	\]
			\item Each of some students owns two computers.
			\item Two computers are owned by each of some students.
			\item Some students possess two computers individually.
			\item Two computers are owned by each of certain students.
			\item $S = Student, \; C = Computer, \; O = Own$
			\item $\exists x (S(x) \wedge \forall y ( C(y) \Rightarrow O(x, y)) )$ \textcolor{red}{\XBox}
			\item $\exists x (S(x) \wedge \exists y, z (( (C(y) \wedge O(x, y) ) \wedge (C(z) \wedge O(x, z) ) ))$ \textcolor{red}{\XBox}
			\item $\exists x (S(x) \wedge \exists y, z (\neg y = z \wedge ( (C(y) \wedge P(x, y) ) \wedge (C(z) \wedge O(x, z) ) ))$ \textcolor{green}{\CheckedBox}
			\item $\exists x (S(x) \wedge \exists y ((C(y) \wedge O(x, twoOf(y)) ))$ \textcolor{green}{\CheckedBox}
		\end{itemize}
	\end{exampleblock}

\end{frame}

\begin{frame}
	\frametitle{\insertshortsubtitle: \insertsection}
	\framesubtitle{\insertsubsection: Propositional meaning (Inference)}

	\begin{exampleblock}{Example of inference}
		\centering
		\begin{tabular}{lll}
			Each human is mortal.  & & Socrates is a human. \\
			$\forall x (Human(x) \Rightarrow Mortal(x))$ && $Human(Socrates)$ \\
			\hline
			\multicolumn{3}{c}{$Mortal(Socrates)$}\\
		\end{tabular}
		
	\end{exampleblock}

\end{frame}

\begin{frame}
	\frametitle{\insertshortsubtitle: \insertsection}
	\framesubtitle{\insertsubsection: Some humor}

	\begin{center}
		\vgraphpage{humor/humor-semantics.jpg}
	\end{center}

\end{frame}

\subsection{Pragmatics and discourse}

\begin{frame}
	\frametitle{\insertshortsubtitle: \insertsection}
	\framesubtitle{\insertsubsection: Pragmatics (Motivation)}
	
	\begin{minipage}{0.7\textwidth}
		\begin{itemize}
			\item Imagine we are at a birthday party;
			\item Someone says: ``Candles!'';
			\item We immediately understand that there are no candles on the cake.
			\item \optword{Question}: how can we infer this?
			\item \optword{Answer}: from the \keyword{context}.
		\end{itemize}
	\end{minipage}
	\begin{minipage}{0.28\textwidth}
		\hgraphpage{pragmatics.pdf}
	\end{minipage}
	
	\vfill
	\begin{tabular}{|p{.4\textwidth}|p{.5\textwidth}|}
		\hline
		\bfseries Semantics & \bfseries Pragmatics \\
		\hline 
		\multicolumn{2}{|l|}{A person \textbf{A} arrived late for an appointment.} \\
		\multicolumn{2}{|l|}{\textbf{B}: ``What time do you think it is?''} \\
		\hline 
		\textbf{A} states the time &  \textbf{A} explains or justifies the delay \\
		\hline
	\end{tabular}
	
	\vfill 
	
	\optword{Example}: \expword{``I'm sorry'' and ``My bad'' may have similar meanings, but their usage depends on context.}
	
\end{frame}

\begin{frame}
	\frametitle{\insertshortsubtitle: \insertsection}
	\framesubtitle{\insertsubsection: Pragmatics (Reasoning)}
	
	\begin{itemize}
		\item \optword{Implicature}: meaning that is implied rather than explicitly stated.
		According to \cite{1979-Grice}, interlocutors are assumed to follow conversational maxims: quantity, quality, relevance, and manner.
		
		E.g., \expword{``Some students passed the exam''} $\Rightarrow$ \expword{``not all students passed''}.
		
		\item \optword{Presupposition}: assumptions taken for granted by the speaker and the listener.
		
		E.g., \expword{``Karim has stopped writing''} $\Rightarrow$ \expword{``Karim used to write''}.
		
		\item \optword{Speech act}: an action performed through language \cite{1962-austin}.
		
		E.g., \expword{declaration, order, question, prohibition, greeting, invitation, congratulation, apology}.
	\end{itemize}
	
\end{frame}

\begin{frame}
	\frametitle{\insertshortsubtitle: \insertsection}
	\framesubtitle{\insertsubsection: Discourse (Coreference)}
	
	\begin{exampleblock}{Example of coreference}
		\begin{itemize}
			\item \keyword{The cat} does not fit in the box because \keyword{it} is too big.
			\item The cat does not fit in \keyword{the box} because \keyword{it} is too small.
		\end{itemize}
	\end{exampleblock}
	
	\keyword{Coreference}: multiple expressions referring to the same entity.
	\begin{itemize}
		\item \optword{Pronouns}: \expword{The cat} chased a mouse and \expword{it} is playing with \expword{it}.
		
		\item \optword{Noun phrases}: \expword{Apple} is a computer manufacturer. \expword{The firm} is known worldwide.
		
		\item \optword{Zero anaphora}: in languages like Japanese, the referent may be omitted.
		
		\item \optword{Named entities / abbreviations}: \expword{International Business Machines (IBM)} is an American company. Like Apple, \expword{IBM} manufactures computers.
	\end{itemize}
	
\end{frame}

\begin{frame}
	\frametitle{\insertshortsubtitle: \insertsection}
	\framesubtitle{\insertsubsection: Discourse (Coherence 1)}
	
	\begin{itemize}
		\item Coherence refers to the logical links between units (sentences or propositions) in a text.
		\item \keyword{Rhetorical Structure Theory (RST)}: a widely used discourse model.
		\item Two types of discourse units: \optword{nucleus} and \optword{satellite}.
		\begin{itemize}
			\item \expword{The student was absent yesterday. He was sick.}
			\item The first sentence is the nucleus (main information).
			\item The second sentence is a satellite (provides an explanation).
		\end{itemize}
		\item Common coherence relations include: reason, elaboration, evidence, attribution, narration, etc.
	\end{itemize}
	
\end{frame}

\begin{frame}
	\frametitle{\insertshortsubtitle: \insertsection}
	\framesubtitle{\insertsubsection: Discourse (Coherence 2)}

	\begin{exampleblock}{Example of ``reason''}
		[\textsubscript{N} The student was absent yesterday.] [\textsubscript{S} He was sick.]
	\end{exampleblock}

	\vfill
	
	\begin{exampleblock}{Example of ``elaboration''}
		[\textsubscript{N} The exam is easy.] [\textsubscript{S} It takes only one hour.]
	\end{exampleblock}

	\vfill
	
	\begin{exampleblock}{Example of ``evidence''}
		[\textsubscript{N} Kevin must be here.] [\textsubscript{S} His car is parked outside.]
	\end{exampleblock}

\end{frame}

\begin{frame}
	\frametitle{\insertshortsubtitle: \insertsection}
	\framesubtitle{\insertsubsection: Some humor}

	\begin{center}
		\vgraphpage{humor/humor-pragmatics.jpg}
	\end{center}

\end{frame}

%===================================================================================
\section{Applications of NLP}
%===================================================================================

\begin{frame}
	\frametitle{\insertshortsubtitle}
	\framesubtitle{\insertsection}
	
	\begin{itemize}
		\item \optword{Tasks}: a task is a subproblem within an NLP application (a building block of a system).
		
		\item \optword{Systems}: a system is a complete application that integrates multiple NLP tasks and is designed for end users (not necessarily familiar with NLP).
		
		\item \optword{Business applications}: NLP can be applied across various domains and industries.
		
		\item Example: Chatbot system = intent classification + entity recognition + dialogue management.
		
	\end{itemize}
	
\end{frame}


\subsection{Tasks}

\begin{frame}
	\frametitle{\insertshortsubtitle: \insertsection}
	\framesubtitle{\insertsubsection: Morphosyntax}

	\begin{itemize}
		\item \optword{Sentence boundary detection and word tokenization}: segmenting text into sentences and words.
		\item \optword{Stemming and lemmatization}: reducing inflected (and sometimes derived) forms to a base form.
		\item \optword{Part-of-speech tagging}: assigning grammatical categories to words in a sentence.
		\item \optword{Parsing}: analyzing the syntactic structure of a sentence.
		\item \optword{Terminology extraction}: identifying domain-specific terms in a text.
	\end{itemize}

\end{frame}

\begin{frame}
	\frametitle{\insertshortsubtitle: \insertsection}
	\framesubtitle{\insertsubsection: Semantics}

	\begin{itemize}
		\item \optword{Word-sense disambiguation}: identifying the correct meaning of a word in context.
		\item \optword{Semantic role labeling}: assigning roles (agent, theme, etc.) to sentence constituents.
		\item \optword{Named Entity Recognition}: identifying entities such as persons, organizations, locations, dates, etc.
		\item \optword{Semantic analysis}: deriving a structured representation of meaning.
		\item \optword{Paraphrase generation}: reformulating a text while preserving its meaning.
		\item \optword{Textual entailment}: determining whether one text logically implies another.
		\item \optword{Automatic text generation}.
	\end{itemize}

\end{frame}

\begin{frame}
	\frametitle{\insertshortsubtitle: \insertsection}
	\framesubtitle{\insertsubsection: Discourse}

	\begin{itemize}
		\item \optword{Coreference resolution}: linking referring expressions that denote the same entity.
		\item \optword{Coherence modeling}: identifying relations between sentences or discourse units.
		\item \optword{Ellipsis resolution}: recovering omitted elements in a sentence.
		Example: \expword{``Karim eats cherries, Mouloud strawberries''}.
	\end{itemize}

\end{frame}

\subsection{Systems}

\begin{frame}
	\frametitle{\insertshortsubtitle: \insertsection}
	\framesubtitle{\insertsubsection: Machine translation}

	\begin{itemize}
		\item Translation of text from one language to another.
	\end{itemize}
	
	\begin{figure}
		\hgraphpage{exp.MT.png}
		\caption{Example of Google Translate.}
	\end{figure}

\end{frame}

\begin{frame}
	\frametitle{\insertshortsubtitle: \insertsection}
	\framesubtitle{\insertsubsection: Automatic text summarization}

	\begin{itemize}
		\item Generating a concise and informative summary of a text.
	\end{itemize}
	
	\begin{figure}
		\hgraphpage{exp.ATS.png}
		\caption{Example of an automatic summarization system.}
	\end{figure}

\end{frame}

\begin{frame}
	\frametitle{\insertshortsubtitle: \insertsection}
	\framesubtitle{\insertsubsection: Question Answering}

	\begin{itemize}
		\item Providing answers to questions expressed in natural language.
	\end{itemize}
	
	\begin{figure}
		\hgraphpage{exp.QR.png}
		\caption{Example of the START question answering system.}
	\end{figure}

\end{frame}

\begin{frame}
	\frametitle{\insertshortsubtitle: \insertsection}
	\framesubtitle{\insertsubsection: Conversational agents (chatbots)}

	\begin{itemize}
		\item Interacting with users through natural language dialogue.
	\end{itemize}
	
	\begin{figure}
		\centering
		\hgraphpage[.5\textwidth]{exp.chatbot.png}
		\caption{Example of Cleverbot.}
	\end{figure}

\end{frame}

\begin{frame}
	\frametitle{\insertshortsubtitle: \insertsection}
	\framesubtitle{\insertsubsection: Information extraction}

	\begin{itemize}
		\item \optword{Text mining}: extracting structured information and knowledge from text.
		\item \optword{Sentiment analysis}: identifying opinions or emotions expressed in text.
		E.g., \expword{determining whether users are satisfied with a product}.
		\item \optword{Recommendation systems}: suggesting items likely to interest a user.
		E.g., \expword{recommending books or products}.
	\end{itemize}

\end{frame}

\begin{frame}
	\frametitle{\insertshortsubtitle: \insertsection}
	\framesubtitle{\insertsubsection: Meme generation}

	\begin{figure}
		\centering
		\vgraphpage[0.7\textheight]{exp.meme.png}
		\caption{Architecture to generate memes \cite{sadasivam2020memebot}.}
	\end{figure}

\end{frame}


\subsection{Business}

\begin{frame}
	\frametitle{\insertshortsubtitle: \insertsection}
	\framesubtitle{\insertsubsection: Commerce}

	\begin{itemize}
		\item \optword{Marketing}: identifying potential audiences for specific products.
		
		\item \optword{Customer service}: using chatbots to answer customer queries, and sentiment analysis to assess customers' opinions on products.
		
		\item \optword{Market intelligence}: monitoring blogs, websites, and social media to analyze trends and competitors.
		
		\item \optword{Recruitment}: filtering CVs to assist in identifying relevant candidates.
	\end{itemize}

\end{frame}

\begin{frame}
	\frametitle{\insertshortsubtitle: \insertsection}
	\framesubtitle{\insertsubsection: E-Governance}

	\begin{itemize}
		\item \optword{Government--citizen communication}: citizens with limited literacy can express opinions via audio or video, which can be transcribed into text; conversely, text can be converted into speech.
		
		\item \optword{Opinion mining}: extracting comments, complaints, and feedback on public policies to assess public sentiment.
	\end{itemize}

\end{frame}

\begin{frame}
	\frametitle{\insertshortsubtitle: \insertsection}
	\framesubtitle{\insertsubsection: Health}

	\begin{itemize}
		\item Structuring medical documents to facilitate their use.
		\item Analyzing large volumes of clinical data.
		\item Supporting diagnosis based on symptoms and medical records.
		\item Generating clinical reports.
		\item Using virtual assistants to monitor and assist patients.
	\end{itemize}

\end{frame}

\begin{frame}
	\frametitle{\insertshortsubtitle: \insertsection}
	\framesubtitle{\insertsubsection: Education}

	\begin{itemize}
		\item Language assessment (reading, writing, and speaking).
		\item Spell checking and grammar correction.
		\item Automatic assessment of students' work (e.g., essays and short answers).
		\item E-learning through conversational agents and interactive systems.
	\end{itemize}

\end{frame}

%===================================================================================
\section{NLP Challenges}
%===================================================================================

\begin{frame}
	\frametitle{\insertshortsubtitle}
	\framesubtitle{\insertsection}
	
	\begin{itemize}
		\item NLP aims to solve problems related to automating human communication.
		\item However, the field faces many challenges.
		\item Many of these challenges are the focus of ongoing research.
	\end{itemize}
	
	
\end{frame}

\subsection{Resources}

\begin{frame}
	\frametitle{\insertshortsubtitle: \insertsection}
	\framesubtitle{\insertsubsection}

	\begin{itemize}
		\item Lack of tools and datasets for low-resource languages.
		\item Manual annotation of training and evaluation corpora.
		\item Processing long documents: limitations in modeling long-range dependencies.
	\end{itemize}
	
\end{frame}

\subsection{Language understanding}

\begin{frame}
	\frametitle{\insertshortsubtitle: \insertsection}
	\framesubtitle{\insertsubsection: Ambiguity and Coreference}

	\begin{itemize}
		\item \keyword{Polysemy}: a word with multiple meanings.
		E.g., \expword{bank: financial institution, river bank}.
		
		\item \keyword{Enantiosemy}: a word with two opposite meanings.
		E.g. \expword{Regret, denotes complaint or nostalgia: ``I regret my childhood''.}
		
		\item \keyword{Trope} (figurative language): 
		metaphors (\expword{``It's raining cats and dogs''}), 
		metonomy (\expword{``The room applauded'' [the people in the room]}), 
		irony (\expword{``What a beautiful day!'' [to mean it is not.]}).
		
		\item Coreference in long texts (ambiguity in reference resolution).
		E.g., \expword{\underline{Table data} is dumped into \underline{a file}, which is sent to \underline{a remote site} where \underline{it} is loaded into the database.}
	\end{itemize}
	

\end{frame}

\begin{frame}
	\frametitle{\insertshortsubtitle: \insertsection}
	\framesubtitle{\insertsubsection: Human aspect}
	
	\begin{itemize}
		\item \optword{Personality}: modeling individual speaking styles and their impact on language use.
		
		\item \optword{Language variation}: non-standard language use (e.g., chat language, Arabizi, Frenglish).
		
		\item \optword{Intention}: sentences may convey implicit meanings beyond their literal interpretation.
		
		\item \optword{Emotions}: meaning can depend on tone and emotional context.
	\end{itemize}
	

\end{frame}


\subsection{Evaluation}

\begin{frame}
	\frametitle{\insertshortsubtitle: \insertsection}
	\framesubtitle{\insertsubsection}
	
	\begin{itemize}
		\item \optword{Manual evaluation}
		\begin{itemize}
			\item Time-consuming.
			\item May require expert annotators.
			\item Subjectivity and inter-annotator agreement issues.
		\end{itemize}
		
		\item \optword{Automatic evaluation}
		\begin{itemize}
			\item Some aspects, such as meaning or coherence, are difficult to evaluate automatically.
		\end{itemize}
	\end{itemize}

\end{frame}

\subsection{Ethics}

\begin{frame}
	\frametitle{\insertshortsubtitle: \insertsection}
	\framesubtitle{\insertsubsection}

	\begin{itemize}
		\item \optword{Privacy}: collected data may contain sensitive personal information, which requires careful handling.
		
		\item \optword{Bias and Discrimination}: training data may introduce or amplify biases in NLP systems.
		
		\item \optword{Misuse}: applications such as surveillance, social engineering, or automation-related job displacement.
	\end{itemize}
	

\end{frame}

\begin{frame}
	\frametitle{\insertshortsubtitle: \insertsection}
	\framesubtitle{\insertsubsection: Some Humor}

	\begin{center}
		\vgraphpage{humor/humor-ethics.jpg}
	\end{center}

\end{frame}

\insertbibliography{NLP01}{*}

\end{document}

